% Part:sets-functions-relations
% Chapter: size-of-sets
% Section: pairing-alt

\documentclass[../../../include/open-logic-section]{subfiles}

\begin{document}

\olfileid{sfr}{siz}{pai-alt}

\olsection{另一种配对函数}

\begin{explain}
除了前述枚举，还有其他枚举 $\Nat^2$ 的方法，使我们更容易求出其逆函数。
这里介绍一种。我们不再用阵列来可视化枚举，而是从一个与（初始为空的）空位相关联的正整数列表开始。
设想按如下方式依次用数对 $\tuple{n,m}$ 填充这些空位。
从第一分量为~$0$ 的数对（即 $\tuple{0,m}$）开始，
将第一个数对（即 $\tuple{0,0}$）填入第一个空位，
然后跳过一个空位，将第二个数对（即 $\tuple{0,1}$）填入下一个空位，
再跳过一个空位，依此类推。
此时，枚举的开头（尚不完整）如下所示：
\[\small
\begin{array}{@{}c c c c c c c c c c c@{}}
\mathbf 1 & \mathbf 2 & \mathbf 3 & \mathbf 4 & \mathbf 5 & \mathbf 6 & \mathbf 7 & \mathbf 8 & \mathbf 9 & \mathbf{10} & \dots \\ \\
\tuple{0,0} &  & \tuple{0,1} &  & \tuple{0,2} &  & \tuple{0,3} & & \tuple{0,4} &  & \dots \\
\end{array}
\]
对仍然空着的位置，用数对 $\tuple{1,m}$ 重复这一过程，
同样每隔一个空位放置一个：
\[\small
\begin{array}{@{}c c c c c c c c c c c@{}}
\mathbf 1 & \mathbf 2 & \mathbf 3 & \mathbf 4 & \mathbf 5 & \mathbf 6 & \mathbf 7 & \mathbf 8 & \mathbf 9 & \mathbf{10} & \dots \\ \\
\tuple{0,0} & \tuple{1,0} & \tuple{0,1} &  & \tuple{0,2} & \tuple{1,1} &
\tuple{0,3} & & \tuple{0,4} &  \tuple{1,2} & \dots \\
\end{array}
\]
以同样的方式填入数对 $\tuple{2,m}$、$\tuple{3,m}$ 等等。
完成后的枚举开头如下：
\[\small
\begin{array}{@{}cc c c c c c c c c c@{}}
\mathbf 1 & \mathbf 2 & \mathbf 3 & \mathbf 4 & \mathbf 5 & \mathbf 6 & \mathbf 7 & \mathbf 8 & \mathbf 9 & \mathbf{10} & \dots \\ \\
\tuple{0,0} & \tuple{1,0} & \tuple{0,1} & \tuple{2,0}  & \tuple{0,2} &
\tuple{1,1} & \tuple{0,3} & \tuple{3,0}  & \tuple{0,4} &  \tuple{1,2} & \dots \\
\end{array}
\]
如果按此枚举为上述阵列中的格子编号，
我们看到的将不是一条整齐的锯齿线，而是这样的排列：
\[
\begin{array}{ c | c | c | c | c | c | c | c }
& \mathbf 0 & \mathbf 1 & \mathbf 2 & \mathbf 3 & \mathbf 4 & \mathbf 5 & \dots \\
\hline
\mathbf 0 & 1 & 3 & 5 & 7 & 9 & 11 & \dots \\
\hline
\mathbf 1 & 2 & 6 & 10 & 14 & 18 & \dots & \dots \\
\hline
\mathbf 2 & 4 & 12 & 20 & 28 & \dots & \dots & \dots \\
\hline
\mathbf 3 & 8 & 24 & 40 & \dots & \dots & \dots & \dots \\
\hline
\mathbf 4 & 16 & 48 & \dots & \dots & \dots & \dots & \dots \\
\hline
\mathbf 5 & 32 & \dots & \dots & \dots & \dots & \dots & \dots \\
\hline
\vdots & \vdots & \vdots & \vdots & \vdots & \vdots & \vdots & \ddots\\
\end{array}
\]
可以看出，第~$0$ 行的数对位于我们枚举的奇数编号位置，
即数对 $\tuple{0,m}$ 在第 $2m+1$ 个位置；
第二行的数对 $\tuple{1,m}$ 位于编号为某个奇数两倍的位置，
具体为 $2 \cdot (2m+1)$；
第三行的数对 $\tuple{2,m}$ 位于编号为某个奇数四倍的位置，
即 $4 \cdot (2m+1)$；以此类推。
各行中 $(2m+1)$ 所乘的因子 $1$、$2$、$4$、$8$、\dots，
恰好是~$2$ 的幂：$1= 2^0$，$2 = 2^1$，$4 = 2^2$，$8 = 2^3$，\dots\@
事实上，相应的指数总是该数对的第一个分量。
因此，对于数对 $\tuple{n,m}$，因子为 $2^n$。
由此得到通式：$2^n \cdot (2m+1)$。
然而，这是将数对映射到\emph{正整数}，
即 $\tuple{0,0}$ 位于位置~$1$。
如果我们要从位置~$0$ 开始计数，
就必须将结果减去~$1$。由此得到：
\end{explain}

\begin{ex}
由
\[
h(n,m) = 2^n (2m+1) - 1
\]
给出的函数 $h\colon \Nat^2 \to \Nat$ 是自然数对集合~$\Nat^2$ 的一个配对函数。
\end{ex}

\begin{explain}
因此，在我们对 $\Nat^2$ 的第二种枚举中，数对 $\tuple{0,0}$ 的代码是 $h(0,0) = 2^0(2\cdot 0+1) - 1 = 0$；
$\tuple{1,2}$ 的代码是 $2^{1} \cdot (2 \cdot 2 + 1) - 1 = 2
\cdot 5 - 1 = 9$；$\tuple{2,6}$ 的代码是 $2^{2} \cdot (2
\cdot 6 + 1) - 1 = 51$。
\end{explain}

有时，只需对自然数对~$\Nat^2$ 进行编码，而不要求该编码是满射。这类编码的逆仅是偏函数。

\begin{ex}
由
\[
j(n,m) = 2^n3^m
\]
给出的函数 $j\colon \Nat^2 \to \Nat^+$ 是一个 $\Nat^2 \to \Nat$ 的单射函数。
\end{ex}

\end{document}